%==============================================================================
% Chapter 4 -- Results and discussion
%==============================================================================
\chapter{Results and discussion}
\label{ch:results}

% -----------------------------------------------------------------------------
% Adapted from Section 4 of the Design-CP paper. Figure/table placeholders below
% point at files you should export from the paper into figures/. Every numerical
% value carries a "% TODO: confirm" where I could not be fully certain of the
% exact figure -- reconcile each against the paper before submission.
% -----------------------------------------------------------------------------

We now show that Design-CP makes it possible to design whole symmetric
nanoparticles end to end, characterise how it scales with the number of GPUs, and
demonstrate that---under strong symmetry---the approach is even feasible on
commodity hardware. We close with a discussion of its limitations.

\section{Design quality and symmetry}
\label{sec:design-quality}

Because context parallelism is an inference-only reparametrisation, it does not
change the model's weights, but it does let us sample systems far larger than
RFD3's native training crop (of order a few hundred tokens and a few thousand
atoms). % TODO: confirm the exact native crop (paper reports ~384 tokens / 5000 atoms).
Operating beyond that crop is out of distribution, and for low-symmetry targets it
can degrade sample quality: a single $\sim$10{,}800-residue monomer generated in
this regime looks unnatural (Figure~\ref{fig:large-designs}). Imposing a
point-group symmetry at the same system size restores protein-like structure,
because---as described in Section~\ref{sec:rfd3-symmetry}---resymmetrisation
collapses the effective degrees of freedom to a single ASU and couples the
subunits, shrinking the design space the model must navigate. Symmetry is thus not
only the biological property we want in a vaccine scaffold but also the very thing
that makes large-assembly generation tractable and well-behaved.

\begin{figure}[t]
  \centering
  % TODO: export Figure 1 from the Design-CP paper into figures/.
  % \includegraphics[width=\linewidth]{fig1_large_designs.pdf}
  \fbox{\parbox[c][4.5cm][c]{0.9\linewidth}{\centering\itshape
    Placeholder -- Design-CP sharding schematic (1D vs 2D) and qualitative
    comparison of large designs with and without an imposed symmetry.}}
  \caption{Symmetric design with Design-CP. \textbf{(a)} the 1D row-sharding and
    2D grid schemes, each square denoting a sub-tensor of the attention matrix
    assigned to one device. \textbf{(b)} qualitative comparison of large designs:
    without symmetry, samples that exceed the native crop degrade; a point-group
    prior mitigates this at the same system size.}
  \label{fig:large-designs}
\end{figure}

\section{Designing icosahedral nanoparticles}
\label{sec:icosahedral}

Using the 2D scheme on a $2\times2$ grid of $95\,\text{GB}$ data-center GPUs,
% TODO: confirm the exact GPU model reported in the paper.
we generated $n{=}40$ icosahedral nanoparticles---$60$ chains of $210$ residues
each, i.e.\ $12{,}600$ residues per particle---together with a paired control of
$n{=}40$ single-chain $210$-residue monomers produced by vanilla single-GPU RFD3.
We benchmark against lumazine synthase (PDB~1NQX), a natural icosahedral particle
used as a vaccine scaffold \parencite{joseph_lumazine_2025,ladenstein_second_2020}.
On per-chain backbone-sanity metrics the Design-CP chains broadly match the vanilla
RFD3 monomers, with a somewhat heavier tail of clash and deviation outliers; the
symmetry-aware interface metrics track the 1NQX baseline (contacts per interface,
minimum inter-chain distances), although a fraction of designs still show steric
overlaps at the interfaces (Figure~\ref{fig:icosahedral}). These results
establish the central feasibility claim of the project: a full icosahedral
assembly can be denoised in a single trajectory, interfaces included, with no
docking step.

\begin{figure}[t]
  \centering
  % TODO: export Figure 2 from the Design-CP paper into figures/.
  \fbox{\parbox[c][4.5cm][c]{0.9\linewidth}{\centering\itshape
    Placeholder -- icosahedral designs vs vanilla RFD3 monomers and the 1NQX
    baseline: per-chain and symmetry-aware interface metrics.}}
  \caption{Designing icosahedral nanoparticles with Design-CP ($n{=}40$,
    $60\times210=12{,}600$ residues), compared against $n{=}40$ single-GPU RFD3
    monomers and the lumazine-synthase reference (PDB~1NQX).}
  \label{fig:icosahedral}
\end{figure}

\section{Scaling: maximum ASU size and inference time vs.\ number of GPUs}
\label{sec:scaling}

To characterise the memory ceiling we swept the ASU length until out-of-memory for
each GPU count $P$, under icosahedral symmetry, for both schemes
(Table~\ref{tab:scaling}). The maximum feasible ASU length grows with the expected
$\sqrt{P}$ trend: distributing the quadratic pair tensor over $P$ devices buys a
linear reduction in per-device pair memory, hence a square-root increase in the
side length that fits. The 2D scheme achieves better wall-clock scaling than 1D
(often on the order of $2\times$ faster at matched size), because its ring
communication overlaps with compute and is restricted to $\sqrt{P}$-sized
subgroups, whereas the 1D scheme pays a global \texttt{ALLGATHER} per block.

\begin{table}[t]
  \centering
  \caption{Scaling of Design-CP under icosahedral symmetry: maximum ASU length
    before OOM and wall-clock time at that maximum, for the 1D and 2D schemes.
    Capacity ratio $=L^{\text{2D}}_{\max}/L^{\text{1D}}_{\max}$; speedup
    $=t^{\text{1D}}/t^{\text{2D}}$.}%
  % TODO: fill values from the paper's Table 1.
  \label{tab:scaling}
  \begin{tabular}{@{}rcccccc@{}}
    \toprule
    $P$ & $L^{\text{1D}}_{\max}$ & $t^{\text{1D}}$ & $L^{\text{2D}}_{\max}$ &
    $t^{\text{2D}}$ & capacity ratio & speedup \\
    \midrule
    1  & -- & -- & -- & -- & -- & -- \\
    4  & -- & -- & -- & -- & -- & -- \\
    9  & -- & -- & -- & -- & -- & -- \\
    16 & -- & -- & -- & -- & -- & -- \\
    \bottomrule
  \end{tabular}
\end{table}

\section{Designing octahedral nanoparticles on commodity GPUs}
\label{sec:octahedral}

Finally, we show that strong symmetry makes large-assembly design feasible even on
workstation-grade hardware. On $16$ NVIDIA RTX~A4000 GPUs ($16\,\text{GB}$ each,
roughly $6\times$ less per-device memory than the data-center parts above) we
generated full octahedral assemblies of $24$ chains of $\sim$178 residues
($\sim$4{,}270 residues total; $n{=}12$).
% TODO: confirm exact residues/chain (paper figures give 176 vs 178).
The designs show healthy backbone-sanity distributions and well-formed
symmetric interfaces (Figure~\ref{fig:octahedral}), illustrating how Design-CP can
be a practical path towards democratising large-assembly protein design rather than
a technique restricted to large clusters.

\begin{figure}[t]
  \centering
  % TODO: export Figure 3 from the Design-CP paper into figures/.
  \fbox{\parbox[c][4.5cm][c]{0.9\linewidth}{\centering\itshape
    Placeholder -- representative octahedral assembly on 16$\times$16\,GB GPUs,
    with backbone-sanity and interface metrics ($n{=}12$).}}
  \caption{Designing octahedral nanoparticles on workstation-grade GPUs: a
    representative $24$-chain assembly generated on $16$ RTX~A4000 ($16\,$GB) GPUs,
    with per-chain and interface metrics.}
  \label{fig:octahedral}
\end{figure}

\section{Discussion}
\label{sec:discussion}

Design-CP reframes symmetric nanoparticle design as the joint denoising of all
atoms of all subunits, and of their inter-subunit interfaces, within a single
generative trajectory---replacing the rigid-body docking of independently designed
oligomers that has defined the field
(Section~\ref{sec:current-methods}). The two schemes trade implementation
complexity against efficiency: 1D is simple and kernel-free, 2D is more involved
but scales better in wall-clock time. Both are model-agnostic in principle---any
generative model with attention and an $\mathcal{O}(n^2)$ pair representation could
adopt them.

The main limitation follows directly from the design choice of keeping the weights
frozen: pushing beyond the native training crop is out of distribution, and for
weakly symmetric or asymmetric targets this causes quality to degrade
(Section~\ref{sec:design-quality}), while even for symmetric designs a fraction of
particles retain steric clashes at their interfaces. Two lines of work follow.
Methodologically, incorporating context parallelism into the \emph{training} loop
(longer-context training) would move the large-assembly regime back in
distribution. Empirically, the in-silico metrics reported here are necessary but
not sufficient: experimental (in-vitro) expression and assembly validation is the
essential next check on whether these designs fold and assemble as intended. The
research programme in Chapter~\ref{ch:future-work} builds on this foundation.
